%fich.tex
%%%%%%%%%%%%%%%%%%%%%%%%%%%%%%%%%%%%%%%%%%%%%%%%%%%%%%%%%%%%%%%%%%%%%%%%
% Fco. Javier Bóhorquez Ogalla
% Descripción del documento:
	
%%%%%%%%%%%%%%%%%%%%%%%%%%%%%%%%%%%%%%%%%%%%%%%%%%%%%%%%%%%%%%%%%%%%%%%%


%%%%%%%%%%%%%%%%%%%%%%%%%%%%%%%%%%%%%%%%%%%%%%%%%%%%%%%%%%%%%%%%%%%%%%%%
%Clase de documento
\documentclass[12pt, spanish]{article}
%%%%%%%%%%%%%%%%%%%%%%%%%%%%%%%%%%%%%%%%%%%%%%%%%%%%%%%%%%%%%%%%%%%%%%%%


%%%%%%%%%%%%%%%%%%%%%%%%%%%%%%%%%%%%%%%%%%%%%%%%%%%%%%%%%%%%%%%%%%%%%%%%	
%Paquetes de lenguaje:
%%%%%%%%%%%%%%%%%%%%%%%%%%%%%%%%%%%%%%%%%%%%%%%%%%%%%%%%%%%%%%%%%%%%%%%%
\usepackage[utf8]{inputenc}
\usepackage[spanish]{babel}
%\usepackage[spanish, activeacute] {babel}	
%\usepackage[spanish]{babel} 				
%\usepackage[latin1]{inputenc}
%%%%%%%%%%%%%%%%%%%%%%%%%%%%%%%%%%%%%%%%%%%%%%%%%%%%%%%%%%%%%%%%%%%%%%%%

%%%%%%%%%%%%%%%%%%%%%%%%%%%%%%%%%%%%%%%%%%%%%%%%%%%%%%%%%%%%%%%%%%%%%%%%
%Gramatic

\usepackage{/usr/share/texlive/texmf-dist/tex/latex/base/bnf}
%%%%%%%%%%%%%%%%%%%%%%%%%%%%%%%%%%%%%%%%%%%%%%%%%%%%%%%%%%%%%%%%%%%%%%%%

%%%%%%%%%%%%%%%%%%%%%%%%%%%%%%%%%%%%%%%%%%%%%%%%%%%%%%%%%%%%%%%%%%%%%%%%
%Paquetes para encabezados y pie de páginas
%%%%%%%%%%%%%%%%%%%%%%%%%%%%%%%%%%%%%%%%%%%%%%%%%%%%%%%%%%%%%%%%%%%%%%%%
\usepackage{fancyhdr}
%%%%%%%%%
% \pagestyle{fancy} %Si se cambia el estilo de página, antes de empezar 
	%el documento. Esto hace que el emcabezado y pie de página se 
	%quede separado del texto por una línea
%%%%%%%%%	
% \fancyhead{} % Límpia el texto que se está usando como encabezado
%%%%%%%%%
% \fancyhead[OPCIONES]{Encabezado}
	% OPCIONES:
		% L texto a la izquierda
		% C texto centrado
		% R texto a la derecha
		% E página par
		% O pagina impar
	%Ejemplo: 
		% fancyhead [LE] {"Doc. Latex"} 
			% Establece como emcabezado de las páginas pares "Doc Latex"
			% con alineacion a la izquierda
%%%%%%%%%			
% \fancyfoot[OPCIONES]{Píe de pag.}
	% OPCIONES:
		% L C R E O
	%Ejemplo
		% \fancyfoot[LE,RO]{\thepage} 
			% Establece como píe de pág. el número de página. con 
			% alineación a la izq en paginas pares y a la derecha en
			% las impares
%%%%%%%%%			
% \renewcommand{\headrulewidth}{0.4pt}
% \renewcommand{\footrulewidth}{0.4pt}
	% Fija el grosor de la la linea que separa el emcabezado y pie de pág
%%%%%%%%%
% Otros comandos:
%\lhead{}
%\chead{}
%\rhead{}
%\lfoot{}
%\cfoot{}
%%%%%%%%%%%%%%%%%%%%%%%%%%%%%%%%%%%%%%%%%%%%%%%%%%%%%%%%%%%%%%%%%%%%%%%%


%%%%%%%%%%%%%%%%%%%%%%%%%%%%%%%%%%%%%%%%%%%%%%%%%%%%%%%%%%%%%%%%%%%%%%%%
% Paquetes para tamaños y distancias
%%%%%%%%%%%%%%%%%%%%%%%%%%%%%%%%%%%%%%%%%%%%%%%%%%%%%%%%%%%%%%%%%%%%%%%%
\usepackage{anysize} 
%%%%%%%%%%%
% \marginsize{3cm}{2cm}{2cm}{2cm}
	% Controla los márgenes {izquierda}{derecha}{arriba}{abajo}. 
%%%%%%%%%%%%%%%%%%%%%%%%%%%%%%%%%%%%%%%%%%%%%%%%%%%%%%%%%%%%%%%%%%%%%%%%


%%%%%%%%%%%%%%%%%%%%%%%%%%%%%%%%%%%%%%%%%%%%%%%%%%%%%%%%%%%%%%%%%%%%%%%%
%Paquetes de carácteres especiales:
%%%%%%%%%%%%%%%%%%%%%%%%%%%%%%%%%%%%%%%%%%%%%%%%%%%%%%%%%%%%%%%%%%%%%%%%
\usepackage{dsfont}	
%Para representar conjuntos matematicos comunes: Z, N, R...
	% mathds{R}, mathds{N},...
%%%%%%%%%%%%%%%%%%%%%%%%%%%%%%%%%%%%%%%%%%%%%%%%%%%%%%%%%%%%%%%%%%%%%%%%

%%%%%%%%%%%%%%%%%%%%%%%%%%%%%%%%%%%%%%%%%%%%%%%%%%%%%%%%%%%%%%%%%%%%%%%%
%Paquetes para insertar gráficos
%%%%%%%%%%%%%%%%%%%%%%%%%%%%%%%%%%%%%%%%%%%%%%%%%%%%%%%%%%%%%%%%%%%%%%%%
\usepackage[dvips]{graphicx}
\DeclareGraphicsExtensions{.pdf,.png,.jpg} %solo para PDFLaTeX
%%%%%%%%%%%%%%%%%%%%%%%%%%%%%%%%%%%%%%%%%%%%%%%%%%%%%%%%%%%%%%%%%%%%%%%%

%%%%%%%%%%%%%%%%%%%%%%%%%%%%%%%%%%%%%%%%%%%%%%%%%%%%%%%%%%%%%%%%%%%%%%%%
%Paquetes de color:
%%%%%%%%%%%%%%%%%%%%%%%%%%%%%%%%%%%%%%%%%%%%%%%%%%%%%%%%%%%%%%%%%%%%%%%%
\usepackage{color}
%%%%%%%%%%%%%%%%%%%%%%%%%%%%%%%%%%%%%%%%%%%%%%%%%%%%%%%%%%%%%%%%%%%%%%%%
%Definición de colores:
\definecolor{gray97}{gray}{.97}
\definecolor{gray75}{gray}{.75}
\definecolor{gray45}{gray}{.45}
%%%%%%%%%%%%%%%%%%%%%%%%%%%%%%%%%%%%%%%%%%%%%%%%%%%%%%%%%%%%%%%%%%%%%%%%
\usepackage{multicol}

%%%%%%%%%%%%%%%%%%%%%%%%%%%%%%%%%%%%%%%%%%%%%%%%%%%%%%%%%%%%%%%%%%%%%%%%
%Paquete de listado de código:
%%%%%%%%%%%%%%%%%%%%%%%%%%%%%%%%%%%%%%%%%%%%%%%%%%%%%%%%%%%%%%%%%%%%%%%%
\usepackage{listings}
%%%%%%%%%%%%%%%%%%%%%%%%%%%%%%%%%%%%%%%%%%%%%%%%%%%%%%%%%%%%%%%%%%%%%%%%
%Configuración del listado:

\lstset { 
	frame=Ltb,
     framerule=0pt,
     aboveskip=0.2cm,
     framextopmargin=0.3pt,
     framexbottommargin=0.2pt,
     framexleftmargin=0.3cm,
     framesep=0pt,
     rulesep=.1pt,
     tabsize=2,
     backgroundcolor=\color{gray97},
     rulesepcolor=\color{black},
     %
     stringstyle=\ttfamily,
     showstringspaces = false,
     basicstyle=\scriptsize\ttfamily,
     commentstyle=\color{gray45},
     keywordstyle=\bfseries,
   	%
     numbers=left,
     numbersep=1pt,
     numberstyle=\tiny,
     numberfirstline = false,
     breaklines=true,
}
%%%%%%%%%%%%%%%%%%%%%%%%%%%%%%%%%%%%%%%%%%%%%%%%%%%%%%%%%%%%%%%%%%%%%%%%


%%%%%%%%%%%%%%%%%%%%%%%%%%%%%%%%%%%%%%%%%%%%%%%%%%%%%%%%%%%%%%%%%%%%%%%%
%Indexado. Índice alfabeticos
\usepackage{makeidx}
\makeindex %para habilitar indices
%%%%%%%%%%%%%%%%%%%%%%%%%%%%%%%%%%%%%%%%%%%%%%%%%%%%%%%%%%%%%%%%%%%%%%%%
%Forma de uso:
% \index{clave} %Para crear un indice de materia
% Supongase que queremos meter en el indice alfabetico referencias a 
% las paginas donde se hace referencia a la clave Producto escalar
% en tal caso en cada una de las páginas pondremos \index{Producto escalar}
% \printindex %Para imprimir el índice alfabetico

% Es necesario una doble compilación del documento. La primera genera un 
% fichero.idx el cual tendremos que procesar con la aplicación makeindex
%tras procesarlo se creará un fichero.ind que contendra el código Latex 
% que se inserta en el documento original. 
% Tras la segunda compilación del documento original se sustituira el
% comando \printindex por el contenido del fichero.ind 

%Podemos cambiar el formato de la clave:
%Ejemplo     				&    	Entrada 		&	Comentario
%\index{hola}				&         hola, 1   	&	Entrada simple
%\index{hola!Pedro}			&		 Pedro, 3 	&	Subentrada bajo ‘hola’
%\index{Juan@\textsl{Juan}}  	&		Juan, 2    	&	Entrada con diseño                        
%\index{Pepa@\textbf{Pepa}}	&		Pepa, 7    	&	Igual que antes
%\index{Loli|textbf}		&		Loli, 3    	&	No de página con diseño                     
%\index{Soraya|textit}		&		Soraya, 5  	&	Igual que antes
%%%%%%%%%%%%%%%%%%%%%%%%%%%%%%%%%%%%%%%%%%%%%%%%%%%%%%%%%%%%%%%%%%%%%%%%


%%%%%%%%%%%%%%%%%%%%%%%%%%%%%%%%%%%%%%%%%%%%%%%%%%%%%%%%%%%%%%%%%%%%%%%%


%%%%%%%%%%%%%%%%%%%%%%%%%%%%%%%%%%%%%%%%%%%%%%%%%%%%%%%%%%%%%%%%%%%%%%%%
%Estilo de página
%%%%%%%%%%%%%%%%%%%%%%%%%%%%%%%%%%%%%%%%%%%%%%%%%%%%%%%%%%%%%%%%%%%%%%%%
%\textwidth 6.75in								  %ancho de texto
%\oddsidemargin -0.2in							%margen izquierdo 
\parskip 0.2in									%espacio parrafos
%%%%%%%%%%%%%%%%%%%%%%%%%%%%%%%%%%%%%%%%%%%%%%%%%%%%%%%%%%%%%%%%%%%%%%%%
\newenvironment{changemargin}[2]{%
\begin{list}{}{%
\setlength{\topsep}{300pt}%
\setlength{\leftmargin}{#1}%
\setlength{\rightmargin}{#2}%
\setlength{\listparindent}{\parindent}%
\setlength{\itemindent}{\parindent}%
\setlength{\parsep}{\parskip}%
}%
\item[]}{\end{list}}
%%%%%%%%%%%%%%%%%%%%%%%%%%%%%%%%%%%%%%%%%%%%%%%%%%%%%%%%%%%%%%%%%%%%%%%%
\usepackage{float}
\restylefloat{table}
\usepackage{placeins}
\usepackage{framed}
\usepackage{epsf}
%%%%%%%%%%%%%%%%%%%%%%%%%%%%%%%%%%%%%%%%%%%%%%%%%%%%%%%%%%%%%%%%%%%%%%%%

\usepackage{titlesec}

\setcounter{secnumdepth}{4}

\titleformat{\paragraph}
{\normalfont\normalsize\bfseries}{\theparagraph}{1em}{}
\titlespacing*{\paragraph}
{0pt}{3.25ex plus 1ex minus .2ex}{1.5ex plus .2ex}


%%%%%%%%%%%%%%%%%%%%%%%%%%%%%%%%%%%%%%%%%%%%%%%%%%%%%%%%%%%%%%%%%%%%%%%%	
%datos del documento
\author{Situación actual\\\\\ Fco. Javier Bohórquez Ogalla}						%autor
\date{}														%fecha

\title{ 
\begin{center}
\includegraphics[scale=0.5]{logo-doc.png}
\end{center} 
}
%%%%%%%%%%%%%%%%%%%%%%%%%%%%%%%%%%%%%%%%%%%%%%%%%%%%%%%%%%%%%%%%%%%%%%%%

%%%%%%%%%%%%%%%%%%%%%%%%%%%%%%%%%%%%%%%%%%%%%%%%%%%%%%%%%%%%%%%%%%%%%%%%
%Documento:
%%%%%%%%%%%%%%%%%%%%%%%%%%%%%%%%%%%%%%%%%%%%%%%%%%%%%%%%%%%%%%%%%%%%%%%%
\lstdefinestyle{nonumbers}{numbers=none}
\begin{document}

\maketitle
\pagebreak
\tableofcontents
\pagebreak
\section {Visión general}
En esta sección se lleva a cabo un análisis de la situación actual en el estudio de la teoría de autómatas y los lenguajes formales. 
Se describirán las necesidades presentes en este campo, y que serán objetivo de las soluciones planteadas. A partir de estos objetivos se llevará 
acabo una descripción de los requisitos que debe cumplir la solución tomada. 

\section{Situación actual}

El estudio de los lenguajes formales es anterior a la concepción de las computadoras. Las matemáticas, la lógica y otras ciencias venían haciendo 
uso de los conceptos de los lenguajes formales para la solución de problemas. 

Una computadora desde un punto de vista teórico es un autómata o máquina de estados, que es capaz de ejecutar una serie de instrucciones descritas
en lenguaje máquina. Por tanto fueron los avances en la rama de la teoría de autómata y los lenguajes formales, entre otros campos, los que permitieron la concepción de 
las primeras computadoras. Actualmente la mayoría de las ingenierías de la información se estudian los conceptos teóricos detrás de las computadoras 
y los lenguajes de programación.

Lo más común es que un estudiante de informática comience sus estudios en estos campos con los autómatas: los tipos que existen, cómo se definen 
y para qué se utilizan. Llegando a comprender conceptos como el de estado, alfabeto, etc. Incluso estudiando definiciones de algunos tipos de lenguajes 
formales como los lenguajes regulares.  

Posteriormente el alumno podría utilizar los conceptos aprendidos para estudiar distintos modelos de computación, los cuales son definidos formalmente y desde un punto de vista teórico.
La máquina de turing y el cálculo lambda son piezas esenciales en este punto del aprendizaje. 

Una vez que se es poseedor del conocimiento base en la teoría de autómatas y los lenguajes formales, el alumno puede aplicar estos conocimientos para estudiar las estructuras, mecanismos 
y demás conceptos detrás de los lenguajes de programación. En este punto se estudian los interpretes y compiladores, y los 
conceptos básicos que hay detrás estos. Conceptos como el léxico, sintaxis y la semántica de los lenguajes de programación, las tablas de símbolos, o las distintos tipos de gramáticas. 

Hasta aquí se habrá obtenido los conocimientos teóricos necesarios y el alumno podría comenzar a dar soluciones prácticas a problemas, aplicando así los conocimientos obtenidos. Así es
común que se comience con el desarrollo de analizadores léxicos y sintácticos sencillos y concretos, para luego aplicarles una semántica. Ejercicios como el desarrollo de 
una calculadora suelen ser habituales. Además se estudian algunas de las herramientas que asisten al proceso del desarrollo de este tipo de tecnologías.

Después del proceso descrito al alumno se le ha brindado la oportunidad de profundizar en un campo con multitud de ramas, técnicas, metodologías y conceptos, que son fruto de años de 
estudio de expertos y apasionados. Podría profundizar en el proceso de compilación o traducción, en las distintas gramáticas, en técnicas de optimización o diseñar sus propias herramientas 
de traducción o interpretación de lenguajes formales.

Por otro lado, en la industria de la tecnología de la información se hace uso de lenguajes muy completos, con gran diversidad de mecanismos y bien consolidados, que son efecto de la evolución 
y las necesidades en el sector. En la mayoría de cursos académicos se estudian estas herramientas desde un punto de vista práctico y de uso. 

La teoría de autómatas y lenguajes formales presenta la base para el estudio de los compiladores e interpretes que son parte fundamental de la industrial 
actual. A pesar de ello no existen herramientas divulgativas, colaborativas e interactivas que, a partir de los conocimientos básicos, ayude a comprender cómo se desarrollan los distintos 
mecanismos y herramientas presentes en la tecnología actual.

Los lenguajes de programación han evolucionado mucho desde que comenzaron ha alejarse del lenguaje máquina. El alumnado actual trabaja sobre los conceptos que le ayudan a entender esta evolución, 
pero no dispone de herramientas o medios para ver cómo estos conceptos son trasladados a un producto real y presente en el día a día de un programador actual. 

\section{Necesidades}
Dada la situación actual se precisa de una herramienta que ayude a comprender cómo se implementa y construye un intérprete para un lenguaje de programación. Es condición necesaria que 
este proceso quede correctamente documentado. La herramienta elaborada deberá ser accesible por cualquier interesado en el tema, que desee profundizar en la práctica del desarrollo
de interpretes y lenguajes de programación. 

Se partirá de los conceptos básicos de la teoría de autómatas y los lenguajes formales, así como de la teoría de compiladores e intérpretes. Se asume pues que el usuario dispone
de este conocimiento.

\section{Objetivos}
Se llegará a construir un lenguaje de programación completo con características presente en la tecnología actual. Este proceso quedará correctamente documentado y se pondrá a disposición pública.
Las características que serán contempladas son:

\begin{itemize}
\item Distintos tipos de datos simples y compuestos.
\item Expresiones lógicas y aritméticas.
\item Expresiones y funciones sobre cadenas.
\item Expresiones y funciones sobre vectores de datos.
\item Expresiones y funciones sobre expresiones regulares.
\item Operadores de conversión de tipos.
\item Mecanismos de entrada/salida.
\item Creación de procesos.
\item Manipulación de ficheros.
\item Funciones de fecha y tiempo.
\item Definición y uso de variables de tipado dinámico.
\item Ámbito de variables.
\item Sentencias de control de flujo condicionales e iterativas.
\item Definición y uso de funciones y procedimientos.
\item Mecanismos y técnicas de la programación funcional.
\item Definición y uso de clases de objetos.
\item Mecanismos y técnicas de la programación orientada a objeto.
\item Integración de módulos que extienden el lenguaje.
\end{itemize}

El interprete desarrollado podrá ser usado de una forma interactiva, permitiendo así la ejecución de
instrucciones bajo demanda. Además deberá tener la capacidad de dar información
estructurada de todo el proceso de interpretación. Otros sistemas podrán utilizar
esta información para ilustrar este proceso de una forma flexible. El intérprete también 
se deberá poder ejecutar en modo servidor, recibiendo así cualquier cadena que hace de código fuente por un puerto TCP e interpretándola.  

Toda la documentación generada a partir del proceso de desarrollo deberá ser estructurada y cumplimentada para 
formar una parte de una aplicación web que la haga accesible. La aplicación web además permitirá el uso online del intérprete, y mostrará información relativa al proceso de interpretación.

El proyecto presentará una licencia de uso libre para que pueda ser usado abiertamente por la comunidad,  a la vez que se nutrirá de las contribuciones de la misma. 

\end{document}
%%%%%%%%%%%%%%%%%%%%%%%%%%%%%%%%%%%%%%%%%%%%%%%%%%%%%%%%%%%%%%%%%%%%%%%%
%Fco. Javier Bohórquez Ogalla
